\begin{thebibliography}{99}
\bibitem{Background}
Shared background of this program, \emph{SUSY positivity research program:
mathematical background}, 11 September 2026, expanded 13 and 17 September 2026,
\texttt{papers/susy-positivity/background.pdf}. Equation numbers of the form
$(1.k)$ below refer to its Section~1.

\bibitem{InverseBulk}
Companion manuscript of this program, \emph{Sphere and Schur sources for the
localized Weil form}, version 0.6,
\texttt{papers/susy-positivity/investigations/inverse-bulk-realization}.

\bibitem{Selection}
Companion manuscript of this program, \emph{A jump-process presentation of the
localized Weil form}, version 0.1,
\texttt{papers/susy-positivity/investigations/source-selection-rules}.

\bibitem{ExclusionMap}
Research note of this program, \emph{The exclusion map: what has been ruled out,
by what argument, with what scope}, 17 September 2026,
\texttt{investigations/inverse-bulk-realization/notes/EXCLUSION\_MAP\_20260917.md}.

\bibitem{Bottleneck}
Research note of this program, \emph{Where the bottleneck actually is, and what
an operator search inside a fixed theory would have to do}, 16 September 2026,
\texttt{brainstorm/inverse-bulk-brainstorm/BOTTLENECK\_AND\_OPERATOR\_SEARCH\_20260916.md}.

\bibitem{GaugeTransfer}
Research note of this program, \emph{Gauge transfer and the first-prime boundary
test}, 14 September 2026,
\texttt{investigations/inverse-bulk-realization/notes/GAUGE\_TRANSFER\_TEST.md}.

\bibitem{Suzuki}
M.~Suzuki, \emph{Weil's quadratic form via the screw function}, 2026,
\href{https://arxiv.org/abs/2606.09096}{arXiv:2606.09096}.

\bibitem{DLMF}
NIST Digital Library of Mathematical Functions, Chapter 5,
especially \S\S5.4, 5.7, 5.9 and 5.11,
\href{https://dlmf.nist.gov/5}{\texttt{dlmf.nist.gov/5}}.

\bibitem{YM2}
A.~Ashtekar, J.~Lewandowski, D.~Marolf, J.~Mour\~ao and T.~Thiemann,
\emph{Closed formula for Wilson loops for $SU(N)$ quantum Yang--Mills theory in
two dimensions}, Journal of Mathematical Physics \textbf{38} (1997), 5453--5482,
\href{https://arxiv.org/abs/hep-th/9605128}{arXiv:hep-th/9605128}.

\bibitem{MM}
Y.~Makeenko and A.~A.~Migdal, \emph{Exact equation for the loop average in
multicolor QCD}, Physics Letters B \textbf{88} (1979), 135--137.

\bibitem{CHMS}
D.~Correa, J.~Henn, J.~Maldacena and A.~Sever, \emph{An exact formula for the
radiation of a moving quark in $\mathcal N=4$ super Yang--Mills}, 2012,
\href{https://arxiv.org/abs/1202.4455}{arXiv:1202.4455}.

\bibitem{GiombiKomatsu}
S.~Giombi and S.~Komatsu, \emph{Exact correlators on the Wilson loop in
$\mathcal N=4$ SYM: localization, defect CFT, and integrability}, 2018,
\href{https://arxiv.org/abs/1802.05201}{arXiv:1802.05201}.

\bibitem{CriticalPath}
Research packet of this program, \emph{Depth extension with an adaptive shift},
10 September 2026,
\texttt{papers/shifted-zeta/storage-depth/}\\
\texttt{\ \ archive/background/1-critical\_path\_research.md}.

\bibitem{CumulativePath}
Research packet of this program, \emph{A depth--shift path controlled by
cumulative storage}, 10 September 2026,
\texttt{papers/shifted-zeta/storage-depth/}\\
\texttt{\ \ archive/background/2-cumulative\_storage\_path.md}.

\bibitem{OWL}
E.~B.~Baker III, \emph{Supersymmetric open Wilson lines}, Journal of High Energy
Physics \textbf{2011} (2011), 140,
\href{https://arxiv.org/abs/1102.4948}{arXiv:1102.4948}.

\bibitem{SuzukiDeBranges}
M.~Suzuki, \emph{On the Hilbert space derived from the Weil distribution},
Canadian Journal of Mathematics, published online 3 November 2025,
\href{https://doi.org/10.4153/S0008414X25101739}{doi:10.4153/S0008414X25101739};
preprint \href{https://arxiv.org/abs/2301.00421}{arXiv:2301.00421}. Read for its
statement of the isomorphism, its structure function and its chain of subspaces;
see Remark~\ref{rem:suzuki}.

\bibitem{BombieriWeil}
E.~Bombieri, \emph{Remarks on Weil's quadratic functional in the theory of prime
numbers, I}, Rendiconti Lincei --- Matematica e Applicazioni \textbf{11} (2000),
183--233. Cited for the attainment of the minimum; read through a scan, and the
negative statements made about it here are "not found" rather than "absent".

\bibitem{OrtegaCerdaSeip}
J.~Ortega-Cerd\`a and K.~Seip, \emph{Fourier frames}, Annals of Mathematics
\textbf{155} (2002), 789--806,
\href{https://arxiv.org/abs/math/0005092}{arXiv:math/0005092}.

\bibitem{OlevskiiUlanovskii}
A.~Olevskii and A.~Ulanovskii, \emph{On Beurling's sampling theorem in
$\mathbb R^n$}, 2011, \href{https://arxiv.org/abs/1106.0576}{arXiv:1106.0576}.
Used for the explicit constant in Beurling's gap theorem.

\bibitem{Kulikov}
A.~Kulikov, \emph{Sharp estimates for eigenvalues of localization operators
before the plunge region}, 2026,
\href{https://arxiv.org/abs/2603.07407}{arXiv:2603.07407}. Preprint. Cited for
the two-sided bound with unequated constants, and for its \S8, which states the
fixed-occupancy limit to be open.

\bibitem{ZhuCompact}
X.~Zhu, \emph{Weil positivity in compact windows: a finite reduction, certified
two-sided bounds, and a Landau--Widom decay law}, 2026,
\href{https://arxiv.org/abs/2608.24827}{arXiv:2608.24827}. Preprint,
unrefereed. \textbf{Its $L$ is half the $L$ used here}; see
Remark~\ref{rem:laws}. Cited for its certified enclosure at its $L=0.8$ and for
its fitted decay law, not for its barrier, which its own Remark 1.5 concedes
bounds only its own method.

\bibitem{ConnesConsaniProlate}
A.~Connes and C.~Consani, \emph{Zeta zeros and prolate wave operators}, Annals of
Functional Analysis \textbf{15} (2024),
\href{https://arxiv.org/abs/2310.18423}{arXiv:2310.18423}. Cited for the
observation that prolate eigenfunctions generate very small eigenvalues of the
Weil quadratic form; no rate is given there.

\bibitem{FrameBoundNote}
Research note of this investigation, \emph{The lower frame bound of the Riemann
zeros on $PW_{L/2}$: a positive infimum, a units correction, and a collapse that
is not arithmetic}, 17 September 2026,
\texttt{notes/FRAME\_BOUND\_AND\_THE\_DENSITY\_20260917.md}. Every table of
Section~\ref{sec:sampling} is reproduced there with its programme.

\bibitem{BhaduriKhareLaw}
R.~K.~Bhaduri, A.~Khare and J.~Law, \emph{Phase of the Riemann zeta function and
the inverted harmonic oscillator}, Physical Review E \textbf{52} (1995),
486--491; and R.~K.~Bhaduri, A.~Khare, J.~Law, M.~V.~N.~Murthy and S.~Tomsovic,
\emph{The phase of the Riemann zeta function}, Pramana \textbf{48} (1997),
537--553. Cited for the phase shift $\arg\Gamma(\frac14+\frac{i\tau}2)$ of the
half-line inverted oscillator, which is \eqref{eq:freewinding}.

\bibitem{MoorePlesserRamgoolam}
G.~Moore, M.~R.~Plesser and S.~Ramgoolam, \emph{Exact S-matrix for
two-dimensional string theory}, Nuclear Physics B \textbf{377} (1992), 143--190,
\href{https://arxiv.org/abs/hep-th/9111035}{arXiv:hep-th/9111035}. Cited only for
the $\Gamma$-structure of the $c=1$ reflection phase used in the survey of
\S\ref{ssec:survey}; a reading, not re-derived here.

\bibitem{Iwaniec}
H.~Iwaniec, \emph{Spectral Methods of Automorphic Forms}, 2nd ed., Graduate
Studies in Mathematics \textbf{53}, American Mathematical Society, 2002. Cited for
the Eisenstein scattering matrix
$\varphi(s)=\Lambda_\zeta(2s-1)/\Lambda_\zeta(2s)$ of the modular surface, with
$\Lambda_\zeta(u)=\pi^{-u/2}\Gamma(\frac u2)\zeta(u)$ --- note this is not the
$\xi$ of \eqref{eq:xi}; the computation in Remark~\ref{rem:modular} is done from
the functional equation directly.

\bibitem{GRT}
S.~Giombi, R.~Roiban and A.~A.~Tseytlin, \emph{Half-BPS Wilson loop and
$AdS_2/CFT_1$}, 2017, \href{https://arxiv.org/abs/1706.00756}{arXiv:1706.00756}.
Cited for the displacement supermultiplet and its protected dimensions.

\bibitem{ZZLiouville}
A.~B.~Zamolodchikov and Al.~B.~Zamolodchikov, \emph{Structure constants and
conformal bootstrap in Liouville field theory}, Nuclear Physics B \textbf{477}
(1996), 577--605,
\href{https://arxiv.org/abs/hep-th/9506136}{arXiv:hep-th/9506136}; and
V.~Fateev, A.~Zamolodchikov and Al.~Zamolodchikov, \emph{Boundary Liouville field
theory I}, 2000, \href{https://arxiv.org/abs/hep-th/0001012}{arXiv:hep-th/0001012}.
The reflection amplitude of Section~\ref{sec:delay} is quoted from the first and
its boundary analogue from the second; neither is re-derived.

\bibitem{ScalingHamiltonian}
A.~Connes and C.~Consani, \emph{The scaling Hamiltonian}, 2019,
\href{https://arxiv.org/abs/1910.14368}{arXiv:1910.14368}. The scaling generator
with continuous spectrum on all of $\R$, which is the object
Section~\ref{sec:delay} identifies as passing the delay clause.

\bibitem{DelayNotes}
Research notes of this investigation, \emph{The delay test, run: no correlator can
wind} and \emph{The delay expansion: the leading term is the degree, the next term
is a selection rule, and it excludes Liouville}, 17 September 2026,
\texttt{notes/THE\_DELAY\_TEST\_20260917.md} and
\texttt{notes/THE\_ARCHIMEDEAN\_EXPANSION\_20260917.md}.

\bibitem{BirmanKrein}
M.~Sh.~Birman and M.~G.~Krein, \emph{On the theory of wave operators and
scattering operators}, Doklady Akademii Nauk SSSR \textbf{144} (1962), 475--478.

\bibitem{Smith}
F.~T.~Smith, \emph{Lifetime matrix in collision theory}, Physical Review
\textbf{118} (1960), 349--356.

\bibitem{SlepianPollak}
D.~Slepian and H.~O.~Pollak, \emph{Prolate spheroidal wave functions, Fourier
analysis and uncertainty --- I}, Bell System Technical Journal \textbf{40}
(1961), 43--63.

\bibitem{BostConnes}
J.-B.~Bost and A.~Connes, \emph{Hecke algebras, type III factors and phase
transitions with spontaneous symmetry breaking in number theory}, Selecta
Mathematica (N.S.) \textbf{1} (1995), 411--457.
\end{thebibliography}
